% ===========================================================================
%  Annexe — Recueil de sujets
% ===========================================================================
\chapter{Recueil de sujets}

\section{La structure de l'épreuve}

\begin{center}
\renewcommand{\arraystretch}{1.4}
\begin{tabular}{|L{4.4cm}|L{7.4cm}|}
\hline
\textbf{Épreuve} & Mathématiques générales\\
\textbf{Série} & Littéraire\\
\textbf{Code matière} & $009$\\
\textbf{Durée} & $2$ heures $15$ minutes\\
\textbf{Matériel autorisé} & Calculatrice scientifique non programmable\\
\textbf{Composition} & Deux exercices et un problème, tous obligatoires\\
\hline
\end{tabular}
\end{center}

\noindent Sous l'ancien programme, le premier exercice portait sur les
suites numériques à chacune des sessions dépouillées pour ce livre, de
$1999$ à $2026$, sans exception. Les suites ne figurant plus au tableau des
contenus du Programme d'Études 2026, ce livre prépare désormais à une
épreuve bâtie sur les trois composantes retenues par le PE :

\begin{center}
\renewcommand{\arraystretch}{1.35}
\begin{tabular}{|L{4.2cm}|c|L{5.6cm}|}
\hline
 & \textbf{Points} & \textbf{Contenu habituel}\\
\hline
Exercice 1 & $5$ & Probabilité\\
Exercice 2 & $5$ & Statistique à deux variables\\
Problème & $10$ & Étude d'une fonction contenant un logarithme ou une
exponentielle ; les dernières questions portent sur une primitive et une
aire\\
\hline
\end{tabular}
\end{center}

\section{Extraits du sujet 2026}

\begin{notepourlaclasse}
La session $2026$ a été composée sous l'ancien programme, et son exercice 1
portait sur les suites numériques. Il n'est pas reproduit ici : ce livre ne
prépare pas cette question, hors du tableau des contenus du PE 2026. Les
deux autres parties du sujet, elles, relèvent entièrement du programme
retenu par ce livre et sont reproduites intégralement.
\end{notepourlaclasse}

\begin{devoirsurveille}[Baccalauréat, Madagascar, série A, session 2026 —
exercice de statistique et problème]
Calculatrice scientifique non programmable autorisée.

\vspace{0.7em}
\noindent\textbf{\sffamily EXERCICE} \hfill \textup{(05 points)}

Le tableau ci-dessous consigne l'évolution de l'effectif des élèves inscrits
aux candidats de l'examen baccalauréat durant six années consécutives.

\begin{center}\footnotesize
\renewcommand{\arraystretch}{1.25}
\begin{tabular}{|l|c|c|c|c|c|c|}
\hline
ANNÉE & $2019$ & $2020$ & $2021$ & $2022$ & $2023$ & $2024$\\
\hline
Rang de l'année $x_{i}$ & $1$ & $2$ & $3$ & $4$ & $5$ & $6$\\
\hline
Nombre d'élèves $y_{i}$ & $340$ & $310$ & $325$ & $380$ & $365$ & $395$\\
\hline
\end{tabular}
\end{center}

\begin{enumerate}[label=\textbf{\arabic*.}, leftmargin=*, itemsep=0.25em]
  \item Construire le nuage des points associé à cette série statistique
        dans un repère orthogonal d'échelle : sur l'axe d'abscisses, $1$ cm
        pour $1$ unité ; sur l'axe des ordonnées, placer $300$ à l'origine
        puis $1$ cm pour $10$ élèves. \bareme{1}
  \item Déterminer les coordonnées du point moyen $G$. \bareme{1}
  \item On partage cette série en deux sous-groupes d'effectifs égaux. La
        première partie est composée des trois premières années et la
        seconde des trois dernières années.
    \begin{enumerate}[label=\textbf{\alph*)}, leftmargin=1.4em, itemsep=0.15em]
      \item Calculer les coordonnées respectives des points moyens $G_{1}$
            et $G_{2}$ de ces deux sous-groupes. \bareme{1}
      \item Établir l'équation de la droite d'ajustement
            $\left(G_{1}G_{2}\right)$ par la méthode de Mayer.
            \bareme{1}
    \end{enumerate}
  \item En utilisant cette droite, estimer l'effectif des élèves inscrits au
        baccalauréat en $2030$. \bareme{1}
\end{enumerate}

\vspace{0.7em}
\noindent\textbf{\sffamily PROBLÈME} \hfill \textup{(10 points)}

On considère la fonction numérique $f$ définie sur
$\intervalleo{0}{+\infty}$ par
\[
  f(x) = -\frac1x-2x+\ln(x) .
\]
On désigne par $(\mathcal C)$ la courbe représentative de $f$ dans un repère
orthonormé $\left(O\,;\,\veci\,,\,\vecj\right)$ d'unité graphique $1$ cm.

\begin{enumerate}[label=\textbf{\arabic*.}, leftmargin=*, itemsep=0.25em]
  \item Calculer les limites de $f$ aux bornes de son ensemble de
        définition. \bareme{1}
  \item \begin{enumerate}[label=\textbf{\alph*)}, leftmargin=1.4em, itemsep=0.15em]
      \item Pour tout $x \in \intervalleo{0}{+\infty}$, calculer $f'(x)$ où
            $f'$ est la fonction dérivée de $f$. \bareme{1,5}
      \item Montrer que $f'(x) = \dfrac{-2x^{2}+x+1}{x^{2}}$,
            $\forall\, x>0$. \bareme{0,75}
      \item En déduire le signe de $f'(x)$, pour tout
            $x \in \intervalleo{0}{+\infty}$. \bareme{0,5}
    \end{enumerate}
  \item Dresser le tableau de variation de $f$. \bareme{1,5}
  \item Écrire une équation de la tangente $(T)$ à $(\mathcal C)$ au point
        d'abscisse $x_{0} = 1$. \bareme{0,75}
  \item Tracer $(\mathcal C)$ et $(T)$ sur l'intervalle
        $\intervalleof{0}{4}$. \bareme{1,5}
\end{enumerate}

\begin{enumerate}[label=\textbf{\arabic*.}, leftmargin=*, itemsep=0.25em, start=6]
  \item Soit $F$ la fonction définie sur $\intervalleo{0}{+\infty}$ par
        $F(x) = -\ln x-x^{2}+x\left(\ln x-1\right)$.
    \begin{enumerate}[label=\textbf{\alph*)}, leftmargin=1.4em, itemsep=0.15em]
      \item Montrer que $F$ est une primitive de $f$ sur
            $\intervalleo{0}{+\infty}$. \bareme{1}
      \item Calculer, en $\text{cm}^{2}$, l'aire du domaine plan délimité
            par la courbe $(\mathcal C)$, l'axe des abscisses et les droites
            d'équations respectives $x = \frac1e$ et $x = 1$.
            \bareme{1,5}
    \end{enumerate}
\end{enumerate}

\vspace{0.4em}
\noindent\emph{On donne $e = 2{,}7$ et $e^{2} = 7{,}4$.}
\end{devoirsurveille}

\begin{notepourlaclasse}
Le corrigé de ces deux parties est réparti dans le livre, à sa place
naturelle : l'exercice au chapitre 12, le problème au chapitre 7 et sa
dernière question au chapitre 9. C'est voulu : un élève qui cherche la
correction d'une question trouve, du même coup, le cours dont elle relève.
\end{notepourlaclasse}

\section{Où retrouver chaque notion dans les sujets}

Le tableau ci-dessous indique, pour chaque notion du programme, les sessions
dont les énoncés figurent dans ce livre et le chapitre où les trouver.

\begin{center}\footnotesize
\renewcommand{\arraystretch}{1.4}
\begin{tabular}{|L{4.1cm}|c|L{6.1cm}|}
\hline
\textbf{Notion} & \textbf{Ch.} & \textbf{Sessions présentes dans le livre}\\
\hline
Second degré, signe d'un trinôme & 1 & $2020$, $2026$ ; bac blanc LPT
$2023$\\
Systèmes linéaires & 2 & $2010$, $2014$, $2017$, $2026$ (extraits de
l'exercice de statistique)\\
Lecture graphique & 3 & $2013$, $2016$, $2017$ (extraits)\\
Limites, asymptotes & 4 & Cameroun probatoire D $2022$, D et TI $2023$ ;
bac blanc LPT $2023$ ; Côte d'Ivoire C $2021$\\
Dérivation, variations & 5 & Cameroun probatoire D $2022$, D et TI $2023$,
IH $2023$ ; bac blanc LPT $2023$\\
Étude complète & 6 & bac blanc LPT $2023$ ; Cameroun probatoire D $2022$,
D et TI $2023$\\
Logarithme népérien & 7 & $2000$, $2002$, $2004$, $2010$, $2012$, $2015$,
$2016$, $2020$, $2026$\\
Exponentielle & 8 & $2001$, $2003$, $2005$, $2013$, $2014$, $2017$,
$2021$ ; bac blanc LPT $2023$\\
Primitives, aires & 9 & $1999$, $2003$, $2005$, $2013$, $2015$,
$2016$, $2017$, $2020$, $2021$, $2026$\\
Dénombrement & 10 & $2000$, $2001$, $2004$, $2009$, $2015$, $2020$,
$2022$\\
Probabilité & 11 & $2004$, $2012$, $2015$, $2016$, $2020$, $2021$,
$2022$\\
Statistique à deux variables & 12 & $2000$, $2002$, $2010$, $2014$,
$2017$, $2026$\\
\hline
\end{tabular}
\end{center}

\noindent Les suites numériques, hors programme du PE 2026, ne figurent plus
dans ce tableau ; les énoncés qui les concernent n'ont pas été retenus pour
ce livre.

\section{Gérer les deux heures quinze}

Deux heures quinze, c'est $135$ minutes pour vingt points : environ sept
minutes par point. Le découpage qui donne les meilleurs résultats est le
suivant.

\begin{center}
\renewcommand{\arraystretch}{1.4}
\begin{tabular}{|L{3.2cm}|c|L{6.4cm}|}
\hline
\textbf{Moment} & \textbf{Durée} & \textbf{Ce que l'on fait}\\
\hline
Lecture & $5$ min & Lire le sujet en entier, repérer la fonction du
problème et les deux thèmes des exercices\\
Exercice 1 & $30$ min & Probabilité : dénombrer avant de calculer,
factorielles et fractions irréductibles\\
Exercice 2 & $30$ min & Statistique : commencer par le tracé du nuage — il
sèche pendant que vous calculez le reste\\
Problème & $60$ min & Dans l'ordre des questions, sans en sauter : chacune
prépare la suivante\\
Relecture & $10$ min & Vérifier les phrases de conclusion, les unités, les
fractions irréductibles\\
\hline
\end{tabular}
\end{center}

\subsection*{Sept conseils qui valent des points}

\begin{enumerate}[label=\textbf{\arabic*.}, leftmargin=*, itemsep=0.25em]
  \item \textbf{Commencez par les exercices, jamais par le problème.} Ils
        valent la moitié du sujet et sont plus courts. Les candidats qui
        échouent sont presque toujours ceux qui ont passé une heure trente
        sur le problème.
  \item \textbf{Une question sautée n'est pas une question perdue.} Les
        questions sont largement indépendantes : si le calcul de $f'(x)$
        résiste, admettez le résultat donné par l'énoncé et poursuivez le
        tableau de variation, qui vaut souvent davantage.
  \item \textbf{Écrivez la phrase de conclusion.} « Donc la droite
        d'équation $x = 0$ est asymptote verticale à $(\mathcal C)$ » vaut
        un demi-point que le calcul seul ne rapporte pas.
  \item \textbf{Respectez l'échelle imposée} pour les tracés, et marquez la
        rupture d'axe quand l'origine des ordonnées est décalée.
  \item \textbf{Donnez les probabilités en fractions irréductibles}, comme
        l'énoncé le demande à chaque session.
  \item \textbf{Utilisez les valeurs approchées fournies.} Quand le sujet
        écrit « on donne $e = 2{,}7$ », c'est qu'il attend un résultat
        numérique : le laisser sous forme littérale coûte des points.
  \item \textbf{Gardez dix minutes.} Pas pour chercher une question de
        plus : pour relire les unités, les signes et les conclusions. C'est
        le quart d'heure le plus rentable de l'épreuve.
\end{enumerate}
